% !TEX root = ../../mixed-sklyanin-cubics.tex
% Current result: 2026-11-22 / v58.
\section{異なる楕円因子を許す全係数の可換比較}
\label{emc:sec:elliptic-comparison}

前節で決定した中心持上げの許容部分に対し、実際の可換比較核を構成する。
各因子の楕円曲線を固定すれば、因子間の $j$ 不変量の一致・相異、
非同種性、非 CM 性、移動点の非捩れ性は必要ない。
二捩れ因子では関係式の線形補正を平行移動で除き、
中心三次自体の三つの一次方向を含める。
本節の主張は、一つの中心変数を共有する因子別 PBW 表示についてのものである。
任意の内在的非可換三次がこの表示を持つことや、楕円因子の moduli 上での
比較の大域的な選択を主張するものではない。

\subsection{固定因子と中心持上げの普遍係数}

$k\ge2$、$N=3k$ とし、各 $i$ について滑らかな点曲線を持つ
標準 Sklyanin 代数
\[
 S_i=S(A_i,B_i,C_i),\qquad
 r_{i,x}=A_i y_i z_i+B_i z_i y_i+C_i x_i^2
\]
とその巡回二関係を固定する。ここで $A_iB_iC_i\ne0$ である。
前節と同じ実楕円データ
$S_i/(\Omega_{i,0})=B(E_i,\mathcal O_{E_i}(1),t_{p_i})$ を取り、
各 $E_i$ を定める非零の平面三次式 $f_i$ を一つ選ぶ。
中心元 $\Omega_{i,0}$ の生成元と重み $a_i\in\C^\times$ も固定する。
二捩れ因子は $(A_i,B_i,C_i)=(1,1,\tau_i)$ に規格化し、
$\Omega_{i,0}$ は前節の $\Omega_{\tau_i}|_{u=0}$ とする。
他の生成元の選択は $a_i$ の変更に吸収できる。

各因子に三つの重み一の係数 $\mathbf t_i$ と、三つの重み二の係数
$\alpha^{(i)}=(\alpha_{i,x},\alpha_{i,y},\alpha_{i,z})$ を置く。
\begin{itemize}
\item $2p_i\ne0$ の場合は $\mathbf t_i=(p_i',q_i',r_i')$ と書き、
\[
 \Phi_i=A_i x_i y_i z_i+B_i y_i x_i z_i
       +\frac{C_i}{3}(x_i^3+y_i^3+z_i^3)
       +uP_i-u^2(\alpha^{(i)}\cdot\mathbf x_i),\qquad
 P_i=p_i'x_i y_i+q_i'y_i z_i+r_i'z_i x_i
\]
の cyclic 微分で周囲因子 $\mathcal B_i$ を定める。
$\Omega_i$ は命題\ref{eco:prop:universal-lift}と
系\ref{eco:cor:non-clifford-lift}の普遍中心持上げである。
純 $u^3$ 係数を零に固定して一意化する。
\item $2p_i=0$ の場合は $\mathbf t_i=(\beta_{i,x},\beta_{i,y},\beta_{i,z})$ とし、
\[
 \mathcal B_i=T[u]\langle x_i,y_i,z_i\rangle/
 (\{y_i,z_i\}+\tau_i x_i^2-u^2\alpha_{i,x},\text{巡回式}),\qquad
 \Omega_i=\Omega_{\tau_i}
        +u(\beta_{i,x}x_i^2+\beta_{i,y}y_i^2+\beta_{i,z}z_i^2)
\]
と置く。これは命題\ref{eco:prop:clifford-center-freedom}の全持上げを含む。
\end{itemize}
係数環と中心商を
\begin{equation}\label{emc:eq:family}
 \begin{gathered}
 T=\C[\mathbf t,\alpha,c],\qquad
 \operatorname{wt}(\mathbf t)=1,\quad
 \operatorname{wt}(\alpha)=2,\quad\operatorname{wt}(c)=3,\\
 B=\bigotimes_{T[u],\,i=1}^k\mathcal B_i,\qquad
 w=\sum_i a_i\Omega_i+cu^3,\qquad A=B/(w),\\
 S=A/(u)=\left(\bigotimes_i S_i\right)/(\sum_i a_i\Omega_{i,0}),
 \qquad G=\sum_i a_i f_i
 \end{gathered}
\end{equation}
とする。係数の内部次数は零であり、表示した重みは法線 $u$ の rescaling に対応する。

\begin{proposition}[指定境界を保つ因子別 PBW 表示の全許容部分]
\label{emc:prop:coverage}
因子間は可換、$u$ は中心とし、各 $S_i$ の関係にその因子の生成元に関する $u$ 倍の線形項と
$u^2$ 倍の定数項を加える全ての PBW 表示を考える。
その中で最高次 active 部分が $\sum_i a_i\Omega_{i,0}$ となる
中心三次が存在するものは、生成元の平行移動
$\mathbf x_i\mapsto\mathbf x_i+u\mathbf s_i$ の後に
式\eqref{emc:eq:family}で表される。
\end{proposition}

\begin{proof}
前節の PBW 条件により線形補正行列は対称である。
平行移動の対角係数は $2C_i\mathbf s_i$ なので、
$C_i\ne0$ により対角を一意に消し、非二捩れの場合は上の $P_i$ を得る。
二捩れの場合、中心性方程式の $u$ 一次・純 $i$ ブロック部分へ射影すると、
定理\ref{eco:thm:clifford-obstruction}と同じ三つの障害を得る。
異なるブロックを含む補正の交換子はこの純ブロック成分へ寄与しない。
従って二捩れの残余線形補正は零でなければならない。
逆にこの条件の下で全ての $\Omega_i$ が存在する。

任意の中心三次 $w'$ から $\sum_i a_i\Omega_i$ を引くと $u q_2$ となる。
PBW により $u$ は正則なので、$q_2$ も中心的である。
その $u=0$ の像は $Z(\bigotimes_iS_i)_2$ に属する。
体上の tensor 積の中心は中心の tensor 積であり、
各因子の $Z(S_i)_1$ は零である。
非二捩れ因子では $Z(S_i)_2=0$、二捩れ因子では
$Z(S_i)_2=\langle x_i^2,y_i^2,z_i^2\rangle$ である。
従って $u\beta$ 倍の平方を引けば差は $u^2q_1$ になる。
再び中心性と $Z(\bigotimes_iS_i)_1=0$ を用いると、残りは $cu^3$ に限られる。
非零の $a_i$ へ係数を吸収すれば表示した族を得る。
\end{proof}

% !TEX root = ../../mixed-sklyanin-cubics.tex
% Current result: 2026-11-22 / v58.
\subsection{固定した楕円因子の周期付き境界比較}
\label{ebk:subsec:boundary}

上で固定した実因子 $S_i$ とその点データ
$(E_i,L_i,t_{p_i})$ を用いる。
各 $f_i$ は同じ偏極曲線 $(E_i,L_i)$ の任意の非零方程式であり、
非可換側の $w_i=a_i\Omega_{i,0}$ と可換側の $a_if_i$ を
それぞれ実際の因子方程式として固定する。
以下の指定周期とは、これらの方程式を用いた
graded MF と Orlov window から移した三乗周期をいう。

\begin{lemma}[任意の固定楕円因子に対する実境界核]
\label{ebk:lem:boundary}
$Z=V(G)\subset\mathbb P^{N-1}$ は滑らかであり、
境界の残余線束圏 $\mathcal K_S$ は全 coherent 残余圏に一致する。
両者は smooth proper で、$\C$ 上に実際の dg Morita 核
\[
 E:\mathcal K_S\xrightarrow{\sim}\Ku(Z)
\]
と rotation の比較射
\[
 \eta:E\varrho_S\xrightarrow{\sim}\varrho_ZE
\]
が存在する。$\eta$ を三回合成した同型は、
$\varrho_S^3\simeq[2]$ と $\varrho_Z^3\simeq[2]$ の
上で指定した周期を保つように取れる。
また、ブロック別スカラーから来る
$Z_0=(\mu_3^k)/\Delta\mu_3$ の指定作用も coherent に保つ。

因子ごとの二捩れ、異なる $j$ 不変量、同種関係や CM を許し、
$3\mid k$ の場合も含む。
\end{lemma}

\begin{proof}
各 $f_i$ の三つの偏微分は、対応する affine 三次元空間で原点以外に
共通零点を持たない。全 $a_i$ が非零で変数ブロックが独立なので、
$G$ の affine critical locus も原点だけであり、$Z$ は滑らかである。

各 $S_i$ は二次 Koszul AS-regular 三次元代数であり、
実際の正則中心元 $\Omega_{i,0}$ による商は
$B(E_i,L_i,t_{p_i})$ である。
主稿の定理\paperref{main}{mbnp:thm:product}の有限 MF 外積と
系\paperref{main}{mbnp:cor:coherent}の飽和を適用できる。
その仮定はこれらの実因子の性質であり、
$j$ の相異性や $2p_i\ne0$ を要求しない。
具体的には、まず三つの augmentation shift が生成する残余線束圏について
有限外積と Shamash 分解により積表示を得る。
積の smooth proper 性と有限群の同変化から残余線束圏の
smooth proper 性を得た後に、同系の coherent 飽和を適用する。
従って全 coherent 圏の生成性を外積比較の前提にしていない。
可換側にも同じ外積と window を適用すると、$Y=\prod_iE_i$ に対し
\begin{equation}\label{ebk:eq:equivariant-products}
 \mathcal K_S^{Z_0}\simeq\Perf(Y)
       \simeq\Ku(Z)^{Z_0}
\end{equation}
という実 dg Morita 比較を得る。
可換側は主稿の系\paperref{main}{mbrp:cor:comparison-field}の比較である。

ここで $Z_0$ は指定されたブロック別スカラー作用である。
その作用は
\[
 \Gamma_k=\{(\lambda_i)\in\mathbf G_m^k:
                    \lambda_i^3=\lambda_j^3\text{ for all }i,j\}
\]
による graded MF と、共通 grading scalar の自明化から得られる。
この構成は実因子の全自己同値群の分類を使わない。
以下の証明は、任意の実因子を扱う主稿の
定理\paperref{main}{pps:thm:elliptic-existence}の証明と、
定理\paperref{main}{csc:thm:comparison}の同変双加群の構成を、
この境界に適用したものである。

右加群の規約では、因子ごとの内部 shift はそれぞれ
\[
 R_{i,p_i}=T_{\OO_{E_i}}\circ(t_{p_i})_*\circ(-\otimes L_i),
 \qquad
 R_{i,0}=T_{\OO_{E_i}}\circ(-\otimes L_i)
\]
である。主稿の楕円 window の補題
\paperref{main}{lem:nfm-twisted-elliptic-window}と
命題\paperref{main}{mbnp:prop:rotations}は、これらを
実方程式からの指定三乗周期とともに同定する。
$\C$ 上では $p_i$ を表す実際の次数零線束 $M_i$ が存在し、
主稿の定理\paperref{main}{thm:nfm-elliptic-normalized-intertwiners}により
\[
 F_i=-\otimes M_i,\qquad
 \eta_i:F_iR_{i,p_i}\xrightarrow{\sim}R_{i,0}F_i
\]
を取れる。
両辺は可逆な perfect Fourier--Mukai 核なので、その自己準同型環は
$\mathrm{HH}^0(E_i)=\C$ である。
$\eta_i$ の三回の合成と両側の指定周期との差は
$u_i\in\C^\times$ であり、$\eta_i$ を $v_i$ 倍すると
差は $v_i^3u_i$ になる。
$v_i^3=u_i^{-1}$ と選んで指定周期を合わせる。
これは主稿の補題\paperref{main}{mbrc:lem:factor-correction}の実構成であり、
自然な方程式直線比較の非消滅を仮定しない。
特に二捩れでその自然な比較が零になる場合にも、
ここで選んだ非零方程式 $f_i$ をそのまま使える。

基準因子 $b$ を一つ選ぶと、式\eqref{ebk:eq:equivariant-products}の
双対指標作用は
\[
 Q_{i,p}=R_{i,p_i}^{-1}R_{b,p_b},\qquad
 Q_{i,0}=R_{i,0}^{-1}R_{b,0}\qquad(i\ne b)
\]
で与えられる。
$F=\boxtimes_iF_i$ と上の $\eta_i$ は
$FQ_{i,p}\simeq Q_{i,0}F$ を定める。
異なる因子の交換は外積の交換則によって保たれ、
$Q_i^3\simeq[-2][2]\simeq\mathrm{id}$ の図式は
合わせた指定周期によって保たれる。
比較射の次数は零で、周期の shift は偶数なので、交換の符号も一致する。
双対群の元を $\prod_{i\ne b}Q_i^{d_i}$、$0\le d_i<3$ の正規形で書けば、
その積は交換と三乗の消去で定まる。
元の MF の中心周期の coherence によりこれらの操作の順序は結果を変えず、
$F$ は全群法則の coherence を intertwine する。
ここで使うのは全次数格子
\[
 \Lambda_k=\mathbb Z^k/\langle3e_i-3e_j\rangle
\]
であり、対称な rotation 分裂を選ばない。
従って $3\mid k$ の場合も同じ証明が成立する。

この導来同変核を実 dg 核として実現する際は、
元の $Z_0$-同変 graded MF 模型へ戻る。
両模型上の $\widehat Z_0$-作用は線型化の character twist による
厳密な dg 作用であり、その間の perfect 双加群の dg 圏にも
両側の twist による厳密な作用がある。
式\eqref{ebk:eq:equivariant-products}で Morita 移送した $F$ は
この双加群圏の対象で、上で確認した図式が
その導来圏での同変構造を与える。
従って\cite[Lemma~8.6]{MEDElagin}を
この\emph{既に厳密な dg 作用を持つ双加群圏}に適用できる。
有限群の誘導、平均化と冪等完備化により、同変構造を
実際の同変 perfect 双加群として取れる。
同補題だけから任意の三角圏作用の dg lift を導いているのではない。
忘却した核は Morita 同値なので、得られた同変核も同値を与える。
\cite[Theorem~4.2, Theorems~8.7 and~8.9]{MEDElagin}の
二重同変化とその dg 実現により、$Z_0$-作用を保つ
$E:\mathcal K_S\simeq\Ku(Z)$ を得る。

最後に、基準因子の $R_{b,p_b}$ と $R_{b,0}$ は、
元の rotation の線型化を積表示へ移したものである。
$FR_{b,p_b}\simeq R_{b,0}F$ の射は、同じ因子比較と外積交換から作られ、
各 $Q_i$ の同変図式と可換する。
従ってこの射も同変双加群圏の射として持ち上がり、
二重同変化後に $\eta$ を与える。
三回の合成の図式は既に合わせた周期の図式であるため、
指定周期も保たれる。
\end{proof}

\begin{remark}[後続の comonad 比較への適用範囲]
\label{ebk:rem:scope}
補題\ref{ebk:lem:boundary}は固定した複素数体値の実因子についての存在結果である。
一度選んだ $E$ は、これらの因子を固定して変形係数だけを動かす
任意の係数環へ定数族として底変更できる。
この $E$ による $K\mapsto EKE^{-1}$ は endofunctor の合成を保つ
Morita 移送なので、実随伴の counit、余乗法と全高次 coherence を運ぶ。
rotation 比較により $\varrho_S^2[-1]$ も
$\varrho_Z^2[-1]$ へ同定される。
これが後続の固定 counit の arity 計算に必要な境界入力であり、
Heisenberg 同変性や異なる因子間で共通な scalar 正規化は不要である。
因子自体を動かす楕円モジュライ上では、実 Picard 線束と周期の三乗根を
族として選ぶ問題が別に残る。
\end{remark}


\subsection{任意の固定因子に対する相対線束模型}

\begin{proposition}[全係数上の周囲正則性と実球面随伴]
\label{emc:prop:line-model}
式\eqref{emc:eq:family}の $B,A$ は左右 noetherian で、各内部次数は有限自由
$T$-加群である。$u$ は $B,A$ 上、$w$ は $B$ 上で中心正則元である。
$B$ は相対的に Koszul で、各複素ファイバーは global dimension と
Gorenstein parameter がともに $N+1$ の AS-regular 代数である。
相対 augmentation 双対は
\[
 R\underline{\operatorname{Hom}}_B(T,B)=T(N+1)[-N-1],\qquad
 R\underline{\operatorname{Hom}}_A(T,A)=T(N-2)[-N]
\]
である。全 twist の線束 dg 圏 $\mathcal C_A$ と、$N-2$ 本の例外線束を
除いた残余 $\mathcal A_A$ は $T$-proper な有限生成模型を持つ。
中心超平面制限 $F:\mathcal A_A\to\mathcal B_T$ とその左右随伴は
全 $T$ 上に存在し、その実随伴の全構造が任意の導来底変更に両立する。
ここで $\mathcal B=\mathcal K_S\simeq\Ku(V(G))$ である。
twist と unshifted cotwist、および生成性は
\[
 \begin{gathered}
 T_F\simeq\varrho^2,\quad C_F\simeq\varrho_A^{-1}[1],\quad
 T_F^3\simeq[4],\quad C_F^3\simeq[1],\\
 \mathcal A_A=\operatorname{thick}(F^R\mathcal B_T),\qquad
 \mathcal B_T=\operatorname{thick}(F\mathcal A_A)
 \end{gathered}
\]
となる。同じ結論は可換族 $G+uQ+u^2L+u^3d$ の全ての係数について成立する。
\end{proposition}

\begin{proof}
命題\ref{lcc:prop:universal-line-model}の証明で因子の特定の係数を使う箇所を
順に確認する。active 次数による PBW 濾過の随伴代数は
$\operatorname{gr}B=(\bigotimes_iS_i)\otimes T[u]$ である。
各固定複素係数での PBW 性は前節で確認済みであり、関係の持上げと
固定した各斉次 $S_i$ の基底により普遍係数上の基底が得られる。
実際、基底の線型関係があれば全複素点で係数が零となり、
多項式環 $T$ の被約性から関係は零である。
各 $S_i$ は中心正則元 $\Omega_{i,0}$ の多項式環上自由である。
これは $S_i/(\Omega_{i,0})$ の斉次基底を持ち上げ、次数に関する帰納法を
使えばよい。従って $\sum_i a_i\Omega_{i,0}$ は tensor 積上正則であり、
\[
 \operatorname{gr}A=S\otimes T[u]
\]
となる。これで次数別自由性と $u,w$ の正則性が従う。
主稿の補題\paperref{main}{mbnp:lem:noetherian}は同一性を仮定しない
楕円 Sklyanin 因子の tensor 積に適用される。
$B/(u)$ の noetherian 性から正次数中心元の持上げで $B,A$ の noetherian 性を得る。

各 $\mathcal B_i$ は $Q_0=T[u]$ 上で cyclic potential の三微分により定義される。
二捩れ因子の $\beta$ は中心元だけを変え、周囲代数の potential には現れない。
cyclic Hessian による長さ三の potential 双加群複体は $Q_0$ 上で定義され、
左右因子交換を伴う転置によって $(3)[-3]$ 自己双対を持つ。
随伴複体は固定した $S_i$ の exact な Koszul 双加群複体であるから、
内部次数ごとに有限な濾過により元の複体も exact である。
ここで標準 $S_i$ の untwisted 性は cyclic potential と
\cite[Corollary~2.9]{MoriSmith}、低次補正後の自己双対性は
\cite[Theorem~3.6, Lemma~3.7 and proof]{NFMBeTa}の同じ明示複体による。
$Q_0$ 上で tensor した長さ $N$ の複体も exact である。

$C=B\otimes_TB^{\mathrm{op}}$ と置く。
共有中心差 $u_L-u_R$ は $C$ 上正則であり、
$D=C/(u_L-u_R)=B\otimes_{Q_0}B^{\mathrm{op}}$ は $C$-perfect である。
$R\underline{\operatorname{Hom}}_C(D,C)=D(1)[-1]$ と上の自己双対から
\[
 R\underline{\operatorname{Hom}}_C(B,C)=B(N+1)[-N-1]
\]
を得る。中心差を因子の個数だけ重複して加えてはいない。
体値の有限濾過 bar 複体の spectral sequence により Koszul 性を得る。
Hilbert 級数は $(1-t)^{-N-1}$ なので Koszul 空間の階数は
$\binom{N+1}{j}$ である。普遍係数上でも各内部次数の行列の階数は
全複素点で一定であり、有限射影の Koszul 空間と底変更に両立する解消を得る。
上の自己双対を片側へ augmentation し、正則中心三次で change of rings
を行えば表示した双対式を得る。最上 Ext の非消滅は global dimension が
正確に $N+1$ であることも示す。

これで命題\ref{lcc:prop:universal-line-model}の局所 cohomology、
線束の Ext 双対、有限 mutation の証明を全て適用できる。
具体的には零次の $A_{l-i}$ と $(N-1)$ 次の
$(A_{i-l-(N-2)})^\vee$ 以外の線束間 Ext は零であり、
$N+1$ 本の連続する線束が全 twist を生成する。
双対に Nakayama twist はなく、制限の二項解消と例外列の長さから
命題\ref{cmm:prop:universal-spherical}と同じ随伴射・周期を得る。
proper 性による高い shift 間の射の消滅が両側の生成性を与える。
全て有限 perfect 模型の構成なので、導来底変更に両立する。
可換側は周囲の多項式環で同じ議論を行える。
非正則な底での全 coherent 圏との同一視はこの段階では主張しない。
\end{proof}

% !TEX root = ../../mixed-sklyanin-cubics.tex
% Current result: 2026-11-22 / v58.
\subsection{一般楕円因子の二つの Bockstein と係数の検出}
\label{ebc:subsec:bockstein}

本小節では、定数補正の二項類と、線形補正または中心平方補正の一次類を検出する。
因子間の同型や Heisenberg 同変な境界比較は仮定しない。
以下で固定する境界は
\[
 B_0=\bigotimes_{i=1}^kS_i,\qquad
 w_0=\sum_i a_i\Omega_{i,0},\qquad
 S=B_0/(w_0),\qquad F=B_0^!
\]
である。中心重み $a_i\ne0$ と区別するため、各因子の Sklyanin パラメータを
$S_i=S(A_i,B_i,C_i)$ と書く。
滑らかな標準点曲線という仮定は $A_iB_iC_i\ne0$ を含む。
$I_{\mathrm{n}}$ を二捩れでない因子の集合、$I_2$ を二捩れ因子の集合とし、
後者は $S_i=S(1,1,\tau_i)$ に規格化する。
前者の一次パラメータは対角零の対称 PBW 行列 $M_i$、
後者の一次パラメータは中心持上げの
$u\sum_j\beta_{ij}x_{ij}^2$ である。
これらをまとめた $\mathbf t$ は $3k$ 次元である。
定数項は potential の $-u^2\sum_{i,j}\alpha_{ij}x_{ij}$ という規約を用いる。

\subsubsection{二次双対と平行移動による商}

\begin{lemma}[一般 Sklyanin 因子の平行移動と独立な左右 gauge]
\label{ebc:lem:translation-gauge}
一因子を $S(a,b,c)$ と省略し、巡回的に
\[
 r_x=ayz+bzy+cx^2,\qquad
 q_x=r_x^*\in (S^!)_2
\]
と置く。二次双対の積は
\begin{equation}\label{ebc:eq:dual-products}
 \xi_x^2=cq_x,\qquad \xi_y\xi_z=aq_x,\qquad
 \xi_z\xi_y=bq_x
\end{equation}
およびその巡回置換である。特に三つの平方は一次独立である。

関係 $r_i+u\sum_jM_{ij}x_j-u^2\alpha_i$ が PBW となる条件は
$M=M^{\mathsf t}$ であり、$\alpha$ は任意である。
$x_i=X_i+s_i u$ とすると線形係数は $M+T(s)$ となる。ただし
\begin{equation}\label{ebc:eq:translation}
 T(s)=
 \begin{pmatrix}
 2cs_x&(a+b)s_z&(a+b)s_y\\
 (a+b)s_z&2cs_y&(a+b)s_x\\
 (a+b)s_y&(a+b)s_x&2cs_z
 \end{pmatrix}.
\end{equation}
従って $\operatorname{Sym}_3/\operatorname{im}T$ は三次元で、
各類は唯一の対角零代表を持つ。

一方、$h_L=\sum L_i\xi_i$、$h_R=\sum R_i\xi_i$ を独立に選び、
$\xi_j\mapsto h_L\xi_j+\xi_jh_R$ の行列を
行が $(q_x,q_y,q_z)$、列が $(\xi_x,\xi_y,\xi_z)$ となるように書くと
\begin{equation}\label{ebc:eq:independent-gauge}
 D(L,R)=
 \begin{pmatrix}
 c(L_x+R_x)&bL_z+aR_z&aL_y+bR_y\\
 aL_z+bR_z&c(L_y+R_y)&bL_x+aR_x\\
 bL_y+aR_y&aL_x+bR_x&c(L_z+R_z)
 \end{pmatrix}.
\end{equation}
この行列は恒等式
\begin{equation}\label{ebc:eq:symmetric-gauge}
 \frac{D(L,R)+D(L,R)^{\mathsf t}}2=T\bigl((L+R)/2\bigr)
\end{equation}
を満たす。従って、独立な左右 gauge の像と対称行列空間の交わりは
正確に $\operatorname{im}T$ である。この結論は $a=b$ でも成立する。
\end{lemma}

\begin{proof}
$(S^!)_2$ は関係空間 $R=\langle r_x,r_y,r_z\rangle$ の双対であり、
$\xi_i\xi_j$ は語 $x_ix_j$ の係数を評価する。
これが式\eqref{ebc:eq:dual-products}を与える。
平方と関係の pairing の行列は $cI_3$ なので、$c\ne0$ の下で独立である。

次数三の唯一の Koszul overlap は
$\sum_i x_ir_i=\sum_i r_ix_i$ である。
その線形補正から得る次数二の式は
$\sum_{i,j}M_{ij}(x_jx_i-x_ix_j)$ である。
これは $\bigwedge^2V$ に属し、一方、各 $r_i$ の平方係数が $c\ne0$ なので
$R\cap\bigwedge^2V=0$ である。
従ってこの式が $R$ に入るためには零でなければならず、これは $M$ の対称性に等しい。
この条件下では低次の PBW 条件も成立する。
実際、線形・定数補正はそれぞれ二次・一次の cyclic potential の微分であり、
\cite[Theorem~3.1]{NFMBeTa}を適用できる。
平行移動の式は三つの二次関係への代入から従う。
対角を消す唯一の値は $s_i=-M_{ii}/(2c)$ である。
同じ Koszul 双加群複体を使えば、これは次数付き代数としての
$\operatorname{HH}^2(S)_{-1}=\operatorname{Sym}_3/\operatorname{im}T$ の計算でもある。
この affine Hochschild 群と残余圏の自然変換空間は、ここでは同一視していない。

式\eqref{ebc:eq:independent-gauge}は双対積から直接得られる。
その対称部分は式\eqref{ebc:eq:symmetric-gauge}である。
対称な $D(L,R)$ は従って平行移動の像に入り、逆に $D(s,s)=T(s)$ である。
\end{proof}

\subsubsection{実自己交叉の類と Tate 模型}

$\kappa$ を $S/S_{\ge1}=\C$ の安定化とし、
$\mathcal T$ を全ての $\kappa(b)$、$b\in\mathbb Z$ を対象とする
$\mathcal K_S$ の full dg 部分圏とする。
rotation はこの模型では $(1)$ であり、指定周期は $(3)\simeq[2]$ である。
非零性の証明には、全 $\mathcal K_S$ 上の nullhomotopy が
$\mathcal T$ へ制限されることだけを用いる。

\begin{lemma}[次数一・二の法線に対する full Bockstein]
\label{ebc:lem:full-tate}
次数 $d\in\{1,2\}$ の中心正則元 $z$ を持つ平坦な持上げ
$\mathcal R$、$\mathcal R/(z)=S$ を考える。
双加群自己交叉三角の connecting class
\begin{equation}\label{ebc:eq:self-intersection}
 S(-d)[1]\longrightarrow S\otimes_{\mathcal R}^{\mathbf L}S
 \longrightarrow S\xrightarrow{\eta_z}S(-d)[2]
\end{equation}
は、$\mathcal R/(z^2)$ の積の一次変形 cocycle の class に等しい。
その加群への作用を singularity quotient と window に通すと、
実自然変換 $\operatorname{id}\to\varrho_S^{-d}[2]$ を得る。

境界の内部次数付き Ext と安定 Hom は、中央周期作用素 $\vartheta$ の加群として
\begin{equation}\label{ebc:eq:tate-module}
 \underline{\operatorname{Ext}}_S^*(\C,\C)
     =F\otimes\C[\vartheta],\qquad
 H^*(\mathcal T)=F\otimes\C[\vartheta,\vartheta^{-1}].
\end{equation}
ここで $F_j$ の双次数は $(j,-j)$、$\deg\vartheta=(2,-3)$ とする。
実自己交叉の full cochain も、この Tate 局所化へ移送される。
\end{lemma}

\begin{proof}
補題\ref{mix:lem:linear-bockstein-characteristic}の bar 比較は
因子の係数を使っていない。その比較を次数 $d$ を明示して記す。
単位を保つ斉次線型切断 $s:S\to\mathcal R/(z^2)$ を選び、
$s(a)s(b)-s(ab)=z\phi(a,b)$ と書く。
$s(a)$ をさらに $\mathcal R$ へ持ち上げ、右 $\mathcal R$-自由解消
$[\mathcal R(-d)\xrightarrow{z}\mathcal R]$ の左作用を表す。
作用の合成と積の作用との差を $z$ で割った homotopy を $S$ へ還元すると、
二側 bar の $[a|b]$ を $\phi(a,b)$ へ送る cochain になる。
これが式\eqref{ebc:eq:self-intersection}の二拡張の代表である。
従って対象ごとの obstruction の値だけでなく、双加群 class 自体が一致する。

自由複体の圏への移送は
Lowen\cite[Theorems~3.21, 4.15]{CMOLowenCharacteristic}の
full brace cochain で行う。
対象が持上げ複体を持つ場合、零項を零に取り、一項は dg Hom の微分の一次補正、
二項は積の一次補正となる。
持上げ複体がない場合も、零項は自由解消の
$\widetilde d^{\,2}/z\bmod z$ であり、同じ class の評価である。
双加群射による自然変換とその homotopy は dg quotient に降下し、
window へ移送される。ここで $\operatorname{HH}^2(S)_{-d}$ 全体が
残余の自然変換空間と同型であるとは主張しない。

$B_0$ は quadratic Koszul AS-regular である。
その Koszul 解消に正則中心三次 $w_0$ の Shamash 構成を施す。
ここでは中心性、正則性と自由解消だけを使う
\cite[Constructions~4.1.1, 4.3.1, Proposition~4.3.2]{CMOEisenbudPeeva}の構成を用いる。
第 $n$ homotopy は homological degree を $2n-1$ 増し、内部 shift が $3n$ なので、
係数の次数は $3n-(2n-1)=n+1>0$ である。
従って augmentation 後の微分は零であり、
周期作用素の加群として式\eqref{ebc:eq:tate-module}の前式を得る。
これは全高次積が自明という主張ではない。

Shamash の divided-power shift は、$B_0\to S$ の
双加群自己交叉から得る中央 cohomology operator $\vartheta$ を表す。
同じ構成を augmented な一次持上げ上で行う場合も、
$\vartheta$ は全相対 Ext 上の中央自然変換である。
augmentation による dg End から Hom 解消への擬同型はその作用と両立する。
従って中央元の derived localization
\cite[Theorem~5.3, Corollary~5.16]{NFMBCLLocalization}を使うと、
射複体は $\vartheta$-telescope で計算できる。
その零ファイバーは中央三次商の Tate 模型
\paperref{moduli}{lem:nfm-cubic-tate-cut}であり、後式を得る。
同じ derived localization の普遍性
\cite[Proposition~3.17, Theorem~3.18]{NFMBCLLocalization}は
持上げた dg 圏と full cochain を同時に運ぶ。
従って局所化した cohomology の次元だけから自然変換を推測しているのではない。
\end{proof}

\begin{proposition}[定数三方向ずつの二項同型]
\label{ebc:prop:constant-iso}
$\mathbf t=0$ での定数補正は線型同型
\begin{equation}\label{ebc:eq:constant-iso}
 \mathfrak b_2:\C^{3k}_{\alpha}\xrightarrow{\ \sim\ }
 H^0\operatorname{RHom}_{\operatorname{Fun}(\mathcal K_S,\mathcal K_S)}
       (\operatorname{id},\varrho_S)
 \xrightarrow{\ \sim\ }\operatorname{Jac}(G)_1
\end{equation}
を与える。最後の同型には補題\ref{ebk:lem:boundary}の固定した境界比較を用いる。
この class は中心 $u$ 制限の正規化二項 comonad データの $\alpha$ による差である。
全 $\mathbf t$ 上でも、二項データの $\alpha$ に関する線型係数は
同じ可逆な線型写像である。
\end{proposition}

\begin{proof}
$\mathbf t=0$、中央 scalar 項 $c=0$ として $v=u^2$、$\deg v=2$ と置く。
各関係は $r_{ij}-v\alpha_{ij}$ となる。
二捩れでない因子の指定中心持上げは
$\Omega_{i,0}+vC_{i,1}(\alpha_i)$ であり、$C_{i,1}$ は active 一次式である。
二捩れ因子では標準 Clifford 体積元をそのまま使える。
これらの中心元の重み付き和で商を取り、$v$ を中心正則元とする $\mathcal R_\alpha$ を得る。
PBW 性と中心元の最高次正則性が、$\C[v]$ 上の次数別自由性を与える。
関係と中心元の $v$ 係数は $\alpha$ に線形なので、前補題は線型写像
$\mathfrak b_2$ を与える。

この射を $\kappa$ 上で評価する。
Koszul--Shamash 解消の最初の二段を $\mathcal R_\alpha$ へ持ち上げると、
微分の二乗を $v$ で割ったものの二次関係成分は
\begin{equation}\label{ebc:eq:constant-evaluation}
 \pm\sum_{i,j}\alpha_{ij}q_{ij}
   =\pm\sum_{i,j}C_i^{-1}\alpha_{ij}\xi_{ij}^2
       \in F_2.
\end{equation}
ここで符号は二拡張の共通の規約だけによる。
中心三次の $v$ 係数は active 一次式なので augmentation 後には零となり、
この成分には寄与しない。
双次数 $(2,-2)$ に $F_j\vartheta^m$ が現れる条件は
$j+2m=2$、$j+3m=2$ であり、唯一の解は $(j,m)=(2,0)$ である。
従って周期成分との混同もない。
各 $q_{ij}$ は独立で、異なる因子は $F_2$ の異なる直和成分に入る。
式\eqref{ebc:eq:tate-module}には $\vartheta$-torsion がないため、
この非零の評価は安定化後も消えない。
全自然変換が零ならその評価も零なので、$\mathfrak b_2$ は単射である。

補題\ref{ebk:lem:boundary}で可換境界へ移す。
対角 MF の sector 公式
\cite[Theorem~5.39 and proof, pp.~82--84]{NFBFKKernels}において、
$\operatorname{Hom}(\operatorname{id},\varrho)$ の恒等 sector は
$\operatorname{Jac}(G)_1$ である。
非恒等 sector の必要条件は $2=N$ であり、$N=3k\ge6$ では成立しない。
Jacobian イデアルは二次式で生成されるので、標的の次元は $3k$ である。
これで式\eqref{ebc:eq:constant-iso}の全単射性が従う。

二項 comonad データとの同定も実解消上で確認できる。
$\widetilde d^{\,2}=v\widetilde\theta$ と書き、$v=u^2$ へ底変更する。
中心法線の二列解消で
\[
 D=\widetilde d+u(\partial_e-\widetilde\theta e),\qquad
 \xi=\partial_e+\widetilde\theta e
\]
と取ると、法線生成元は $\xi^2=\widetilde\theta$ を満たす。
従ってその二項作用は同じ $v$-Bockstein である。
これは補題\ref{cmo:lem:central-universal-parameters}の実随伴の計算であり、
因子固有の係数を使っていない。

命題\ref{emc:prop:line-model}の全係数上の実構造と、
後の補題\ref{euk:lem:arity}の自由な二項比較群を用いると、比較 class は多項式となる。
二項 class の重みは二で、$\operatorname{wt}(\mathbf t)=1$、
$\operatorname{wt}(\alpha)=2$ である。
従って $\alpha$ に関する係数に正次数の $\mathbf t$ は現れず、
全 $\mathbf t$ 上で同じ線型写像となる。
ここで因子間の共通 scalar や標準座標での対角性は結論していない。
\end{proof}

\subsubsection{一次類の一項成分と Clifford 中心の二項成分}

\begin{lemma}[一般線形 PBW の一項成分]
\label{ebc:lem:unary}
$\epsilon^2=0$、$\deg\epsilon=1$ とし、中心 $u$ 変形を一次で切った
$S_\epsilon$ を考える。$\mathcal T$ 上の実 Bockstein を表す full cocycle
$\delta=(\delta_0,\delta_1,\ldots)$ は、$\delta_0=0$ と取れて、
各因子の次数一生成元上で
\begin{equation}\label{ebc:eq:unary}
                 \delta_1(\xi_{ij})=\pm\sum_l M^{(i)}_{lj}q_{il}
\end{equation}
となる。中心持上げの $uC_2$ および二捩れ因子の $u\beta x^2$ は、
この一項成分へ寄与しない。
\end{lemma}

\begin{proof}
全 active 変数を零へ送る $S_\epsilon\to\C[\epsilon]/(\epsilon^2)$ が存在する。
周囲 PBW 代数の相対 augmented Koszul 解消は、通常の Koszul 微分に
$\epsilon$ と線形補正の双対の積を加えたものである。
その完全性は active 濾過の随伴複体から従う。
これは非斉次二次双対の構成
\cite[\S\S2.1--2.5]{CMOPositselskiQL}を、係数環上で行ったものである。
augmentation 後の微分の $\epsilon$ 係数は、関係の双対基底では
式\eqref{ebc:eq:unary}となる。

中心三次の Shamash homotopy の係数は、前補題と同じく総内部次数 $n+1$ である。
$\epsilon^2=0$ なので、その $\epsilon$ 係数も少なくとも $n\ge1$ 個の
active 変数を持ち、augmentation 後に消える。
従って相対 Hom 解消の微分は $\pm\epsilon\ell^*$ で、
$\ell^*$ の $F_1\to F_2$ 成分が表示式である。
補題\ref{ebc:lem:full-tate}の full cochain と中央周期の局所化を使い、
最小 Tate 模型へ移送する。
一項成分は Hom 複体の一次微分が cohomology に誘導する写像なので、
この計算は移送後も同じである。
\end{proof}

\begin{lemma}[Clifford 中心平方の二項成分]
\label{ebc:lem:clifford-binary}
全ての PBW 線形補正と定数補正を零とし、二捩れ因子の中心だけを
\[
       w_\epsilon=w_0+
             \epsilon\sum_{i\in I_2}a_i\sum_j\beta_{ij}x_{ij}^2
\]
へ変形する。$\mathcal T$ の最小模型と full cocycle を選んで
$\delta_0=\delta_1=0$ とでき、その二項成分は同じ Clifford ブロック内で
\begin{equation}\label{ebc:eq:clifford-binary}
        \delta_2(\xi_{ij},\xi_{il})
                 =\pm a_i\beta_{ij}\delta_{jl}\vartheta
\end{equation}
となる。
さらに零ファイバーの $m_3:F_1^{\otimes3}\to\C\vartheta$ は、
当該ブロックに制限すると完全交代であり、
三つの入力が同じブロックに属さなければ零となる。
\end{lemma}

\begin{proof}
この場合、周囲代数の相対 Koszul 解消は一定である。
前補題の次数計算によって、相対 Shamash 解消 $P_\epsilon$ には
augmentation 後の微分が全く残らない。
従って相対 Ext は係数環上自由で、その次数一基底は $\xi_{ij}$、
次数二基底は二次関係の双対と中心関係の双対 $\vartheta$ である。
dg End から相対 Hom への augmentation 擬同型を使って最小模型を取る。
零ファイバーの contraction を延長する perturbation は
$\epsilon^2=0$ により一次で止まる。
相対 cohomology が上の自由加群であるため、移送した微分も零である。
実際の持上げ複体を使っているので曲率も零であり、$\delta_0=\delta_1=0$ となる。

二項成分は相対 Yoneda 積の低次 bar 計算だけで求められる。
自由な最小関係基底を、周囲の二次関係と $w_\epsilon$ で取る。
次数一の二つの Ext class の積とこの関係基底との pairing は、
各関係の active 語長二の係数である。
具体的には、長さ $n$ の語に対応する bar 比較を
$[x_{j_1}\mid x_{j_2}\cdots x_{j_n}]$ で始めると、
二つの一次 cochain の積は $n=2$ だけを残す。
長さ三以上の項の第二因子は augmentation ideal の平方に属し、一次 cochain が消す。
この計算は自由な関係基底と係数環上の線型 splitting だけを使い、
係数環が体であることを必要としない。
$w_\epsilon$ の語長二成分は
$\epsilon\sum a_i\beta_{ij}x_{ij}^2$ なので、
相対積の $\epsilon$ 係数が式\eqref{ebc:eq:clifford-binary}となる。

零ファイバーの三項だけには、体上の関係と高次積の定理
\cite[Theorem~A, Corollary~B]{NFMLPWZExt}を用いる。
Clifford ブロックでは二次関係が対称であり、$S_{i,2}$ の語代表を
$\bigwedge^2V_i$ と三つの平方の和に取れる。
従って
\[
 \Omega_{i,0}=\frac{\tau_i^3-1}{3}
 (yxz+xzy+zyx-xyz-yzx-zxy)
\]
そのものが $V_i\otimes S_{i,2}\subset V_i^{\otimes3}$ 内の関係代表となる。
他因子でも純ブロックの代表を取れば、$w_0$ はそのような三次関係の和である。
同定した $\operatorname{Ext}^1$、$\operatorname{Ext}^2$ の下で
$m_3$ はこの三次関係の双対に取れるため、主張する完全交代性と混合入力の消滅を得る。
この最小模型の選択は Yoneda 積を変えない。
零ファイバーで必要な高次の模型変更を相対模型へ延長しても、
$m_1=0$ なので先に計算した $\delta_2$ は変わらない。
ここでは同文献の体上の定理を、非被約係数環上の定理として無条件に適用してはいない。

最後に中央 cohomology operator による局所化を
補題\ref{ebc:lem:full-tate}と同じ dg 圏上で行う。
通常 Ext を Tate 模型の非負周期部分として回復する補題
\paperref{moduli}{lem:nfm-ordinary-ext-intrinsic-cut}を用いる。
その次数計算では、$m_n$ の出力の周期指数は入力の指数の和に $n-2$ を加えたもの、
一次補正 $\delta_n$ では $n-1$ を加えたものである。
$m_0=m_1=\delta_0=\delta_1=0$ なので、この部分は相対的にも保たれる。
stabilization の最小模型間の射はこの部分で一次成分が同型である。
より具体的には、相対 Ext と Tate はともに微分零の自由模型であり、
局所化射の一次成分は係数環上の分裂単射
$F[\vartheta]\hookrightarrow F[\vartheta,\vartheta^{-1}]$ である。
この $A_\infty$ 射の高成分を、arity ごとにその像上から全体へ斉次線型に延長し、
一次成分が恒等の target の座標変更で順に消す。
bar coalgebra の長さについてのこの再帰により、
選んだ通常 Ext 模型を Tate 模型へ厳密に包含できる。
微分が零なので、この座標変更は二項 Yoneda 積とその一次補正を変えない。
従ってこれらは通常 Ext だけの式ではなく、$\mathcal T$ 上で実際の
自己交叉自然変換を表す cochain の式である。
\end{proof}

\begin{theorem}[持上げ可能な一次パラメータ全体の単射性]
\label{ebc:thm:first-injection}
本節の $3k$ 個の一次パラメータは線型単射
\begin{equation}\label{ebc:eq:first-injection}
 \mathfrak b_1:\C^{3k}_{\mathbf t}\lhook\joinrel\longrightarrow
 H^0\operatorname{RHom}_{\operatorname{Fun}(\mathcal K_S,\mathcal K_S)}
       (\operatorname{id},\varrho_S^2)
       \simeq\operatorname{Jac}(G)_2
\end{equation}
を与える。定数補正 $\alpha$ と中央 scalar 項 $c$ はこの写像に現れない。
特に中心超平面の一次 counit class が零であることは
$\mathbf t=0$ と同値である。
ここで自然変換の homotopy に rotation 不変性や周期不変性を追加していない。
\end{theorem}

\begin{proof}
$u^2=0$ の関係式と中心式は $\mathbf t$ に線形で、
$\alpha,c$ は含まない。従って補題\ref{ebc:lem:full-tate}が線型写像を与える。
可換境界の対角 MF の公式では、次数二の自然変換の恒等 sector は
$\operatorname{Jac}(G)_2$ であり、非恒等 sector の必要条件 $4=N$ は成立しない。
これが表示した標的の同定である。

この class が零と仮定する。全自然変換の nullhomotopy を
$\mathcal T$ に制限し、全 bar cochain $h=(h_0,h_1,\ldots)$ を使って
$\delta=[m,h]$ と書く。
ここで係数は $\varrho_S^{-1}$ の graph 双加群であり、
各 $h_n$ は各対象列に独立に依存してよい。
その cohomological degree は $1-n$、出力の内部 shift は入力の和より一つ小さい。
全 arity の直積を使っており、bar 長を有限に切っていない。

$F_j\vartheta^m$ の双次数は $(j+2m,-j-3m)$ である。
零項 $h_0$ の条件は
\[
 j+2m=1,\qquad j+3m=1,
\]
従って $(j,m)=(1,0)$ であり、$h_0(b)\in F_1$ である。
負の $m$ を排除してから解いたのではなく、全整数 $m$ に対する唯一の解である。
$m_0=m_1=0$ なので、arity one の境界は $[m_2,h_0]$ だけである。
一つの隣接する対象対で左右の $h_0$ を独立な $h_L,h_R$ とし、
第 $i$ ブロックの $F_{i,2}$ へ射影する。
補題\ref{ebc:lem:unary}の左辺は対称行列 $\pm M_i$、
右辺は式\eqref{ebc:eq:independent-gauge}である。
他ブロックの $h_0$ はこの射影では消える。
補題\ref{ebc:lem:translation-gauge}より $M_i\in\operatorname{im}T_i$ であるが、
対角零の正規化と $C_i\ne0$ によって $M_i=0$ となる。
これで全ての $i\in I_{\mathrm n}$ の一次パラメータが消える。

残る変形は補題\ref{ebc:lem:clifford-binary}の中心平方変形である。
同じ full nullhomotopy の arity two を取ると、可能な項は
$[m_2,h_1]+[m_3,h_0]$ だけである。
$h_1(F_1)$ の双次数条件と $h_1(F_2)$ の条件はそれぞれ
\[
 \begin{array}{c|c|c}
 & (j+2m,j+3m)&(j,m)\\\hline
 h_1(F_1)&(1,2)&(-1,1)\\
 h_1(F_2)&(2,3)&(0,1)
 \end{array}
\]
である。従って $h_1(F_1)=0$、$h_1(F_2)\subset\C\vartheta$ であり、
$[m_2,h_1]$ から残るのは $h_1\circ m_2$ だけである。

一つの Clifford ブロック $S(1,1,\tau)$ と、二つの連続する次数一の射を固定する。
出力の $\C\vartheta$ 成分を見ると、$h_1\circ m_2$ の行列は、共通の符号を吸収して
\begin{equation}\label{ebc:eq:clifford-binary-gauge}
 \begin{pmatrix}
 \tau\lambda_x&\lambda_z&\lambda_y\\
 \lambda_z&\tau\lambda_y&\lambda_x\\
 \lambda_y&\lambda_x&\tau\lambda_z
 \end{pmatrix}
\end{equation}
となる。ここで $\lambda_i$ は $h_1(q_i)$ の係数で、対象対ごとに任意でよい。
一方、$[m_3,h_0]$ の三つの挿入は、完全交代な三次 tensor へ
それぞれ $h_0(b)$、$h_0(b-1)$、$h_0(b-2)$ を挿入するものである。
三つの $h_0$ は互いに異なってよいが、各挿入は残る二入力について交代的である。
入力の生成元を交換しても対象列は同じなので、対象依存性はこの交代性を壊さない。
他ブロックからの挿入は、純ブロックの $m_3$ によって零となる。

従って arity two の式の対称部分は
\[
 \pm a_i\operatorname{diag}(\beta_{ix},\beta_{iy},\beta_{iz})
   =
 \begin{pmatrix}
 \tau_i\lambda_x&\lambda_z&\lambda_y\\
 \lambda_z&\tau_i\lambda_y&\lambda_x\\
 \lambda_y&\lambda_x&\tau_i\lambda_z
 \end{pmatrix}.
\]
対角外成分から全ての $\lambda$ が零となり、$a_i\ne0$ より全ての $\beta_{ij}$ も零となる。
各 Clifford ブロックで同じ議論を行えば $\mathbf t=0$ を得る。
逆に $\mathbf t=0$ なら $S_\epsilon=S[\epsilon]/(\epsilon^2)$ なので一次 class は零である。
\end{proof}

\begin{corollary}[固定楕円因子ごとの一次像]\label{ebc:cor:pure-image}
補題\ref{ebk:lem:boundary}の外積から作った境界核と rotation 比較を用いると、
$\mathfrak b_1$ は各 $\C^3_{\mathbf t_i}$ を
$\operatorname{Jac}(f_i)_2\subset\operatorname{Jac}(G)_2$ に同型に写す。
従ってその像はこれらの直和である。$\mathfrak b_2$ も各
$\C^3_{\alpha^{(i)}}$ を第 $i$ ブロックの一次式全体に同型に写す。
\end{corollary}
\begin{proof}
基準因子 $b$ を固定する。$H_b=\{(s_i)\in\mu_3^k:s_b=1\}$
はブロック別スカラーの射影群 $Z_0$ へ同型に写る。
$x_i\mapsto s_i x_i$（各ブロックの三変数すべて）、$u\mapsto u$ とし、
同じ持上げで両側の内部 shift と可換させる。この線型化が対応する理由は、
境界補題の積模型で内部 shift の持上げが全次数格子の $e_b$、
すなわち $R_b$ だからである。$e_b$ の文字は $H_b$ 上自明であり、
余分な文字 twist はない。同補題の $FR_{b,p_b}\simeq R_{b,0}F$ は
各 $R_i^{-1}R_b$ の比較射と可換する。従って二重同変化して得る
rotation 比較もこの線型化を保つ。三乗周期を定める各方程式は $s_i^3=1$ により不変である。

この代入を変形環の間の同型とみる。非二捩れ因子では二次関係を
$s_i^2$ で割ると、$u$ を伴う一次式の係数は $s_i^{-1}\mathbf t_i$、
定数補正は $s_i^{-2}\alpha^{(i)}$ になる。
指定中心持上げも同じ代入で移る。実際、最高項と純 $u^3$ 係数の規格化は
不変であり、系\ref{eco:cor:non-clifford-lift}の一意性が適用できる。
二捩れ因子では体積元は不変（各項のブロック次数は三、または
$u^2\alpha^{(i)}$ を用いて同じ文字になる）で、中心平方の補正係数は
$s_i^2\mathbf t_i$ になる。従って両型で $\mathbf t_i$ の文字は $s_i^2$、
$\alpha^{(i)}$ の文字は $s_i$ である。
法線 $u$（定数方向では $v=u^2$）を固定する同型は、補題\ref{ebc:lem:full-tate}の
二側 bar の積の差を法線で割った類をそのまま運ぶ。これと上の
線型化から両 Bockstein は $H_b$-同変である。
恒等 sector の多項式乗法に対する作用も同じ代入である。

$\operatorname{Jac}(G)=\bigotimes_i\operatorname{Jac}(f_i)$ の次数二は
各純二次部分（文字 $s_i^2$）と、各異なる二因子の一次部分の積
（文字 $s_i s_j$、$i<j$）に分解する。$s_b=1$、残りの文字は
独立な位数三の生成元なので、純部分の文字は互いに異なり、
いずれも混合部分に現れない。各純部分の次元は
$6-3=3$ である。滑らかな平面三次の三偏微分は独立だからである。
定理\ref{ebc:thm:first-injection}の単射性で各像はその文字空間を尽くす。
次数一の文字は $s_i$ で互いに異なり、命題\ref{ebc:prop:constant-iso}が
定数方向の主張を与える。
\end{proof}

\begin{remark}[検算と結論の範囲]
\label{ebc:rem:scope}
\path{computations/verify_elliptic_bockstein.py} は、
元の関係から双対積を作り、式\eqref{ebc:eq:symmetric-gauge}を多項式恒等式として確認する。
一般の独立左右 gauge 六列に対角外対称三列を加えた行列式は
$-8c^3(a-b)^3$ であり、$a=b$ では gauge の階数が三へ下がる。
対称 gauge と平行移動像の一致自体は、この特殊化でも成立する。
Clifford の二項検査では、関係三列・交代三列・中心対角三列の行列式が $8$ となる。
また、全ての Tate 逆冪を許した双次数連立と、対象ごとの独立な三挿入の交代性を確認する。
これらは有限線型代数の証明書であり、実自己交叉から full cochain への比較、
Tate 局所化、残余圏への移送は本文の証明による。

系\ref{ebc:cor:pure-image}は外積境界比較での像を決定する。各三次元ブロック内の行列は未決定である。
同様に式\eqref{ebc:eq:constant-iso}は可逆な線型写像を与え、
異なる因子の共通 scalar や因子モジュライ上の標準的な正規化を与えるものではない。
後続の実比較の存在そのものと、これらのパラメータ非退化性を区別する。
\end{remark}

% !TEX root = ../../mixed-sklyanin-cubics.tex
% Current result: 2026-11-22 / v58.
\subsection{任意の固定楕円因子の普遍 comonad 比較と実核}
\label{euk:subsec:universal-kernel}

以下では本節の因子、中心元、非零重みと係数環
$T=\C[\mathbf t,\alpha,c]$ を固定し、
補題\ref{ebk:lem:boundary}の周期付き実境界比較を一度選ぶ。
二捩れ因子の $\mathbf t$ は関係の線形補正ではなく、
中心持上げへ加える $u$ 倍の中心平方の係数である。
どちらの場合も重みは
\[
 \operatorname{wt}(\mathbf t)=1,\qquad
 \operatorname{wt}(\alpha)=2,\qquad \operatorname{wt}(c)=3
\]
となる。$N=3k$ 個の active 変数を $\mathbf x=(x_1,\ldots,x_N)$ と並べ直し、
\[
 G=\sum_i a_if_i,\qquad Z=V(G)\subset\mathbb P^{N-1},\qquad
 \mathcal B=\Ku(Z),\qquad V_j=\operatorname{Jac}(G)_j
\]
と書く。各 $f_i$ の平面三次が滑らかであるから、
$\partial G$ の共通零点は原点だけであり、$Z$ は滑らかである。
特に $V_1=\C[\mathbf x]_1$、$V_0=\C$ である。

これから構成する Morita 比較には、一次 counit の係数写像の単射性を用いない。
また二項の定数方向が同型であることは、比較の存在を構成した後、
係数変更が多項式同型であることと判別幾何の帰結に使う。
共通の Heisenberg 表現や、全因子を交換する対称性は要求しない。

\subsubsection{一般の滑らかな三次境界における非分裂比較群}

\begin{lemma}[因子の方程式に依存しない正規化塔]
\label{euk:lem:arity}
$\varrho$ を $\mathcal B$ の rotation とし、$\varrho^3\simeq[2]$ の周期を固定する。
任意の可換 $\C$-代数 $R$ と $\eta\in V_2\otimes R$ に対して
\[
 K_\eta=\operatorname{fib}(\mathbf1\xrightarrow{\eta}\varrho^2),
 \qquad \epsilon:K_\eta\to\mathbf1,
 \qquad J=\operatorname{fib}(\epsilon)=\varrho^2[-1]
\]
と置く。$r\ge2$、$i\ge0$ に対して、自然に
\begin{equation}\label{euk:eq:arity}
 H^{2-r-i}\operatorname{RHom}(K_\eta,J^r)=
 \begin{cases}
  V_1\otimes R,&(r,i)=(2,0),\\
  R,&(r,i)=(3,0),\\
  0,&\text{その他}
 \end{cases}
\end{equation}
となる。これらの同定は任意の係数変更と両立する。
固定 counit 上の正規化 $A_r$-comonad 塔の非空 fiber は離散であり、
$r\ge4$ では可縮である。
また $H^{-j}\operatorname{RHom}(\mathbf1,\varrho^2)=0$ $(j\ge1)$ である。
\end{lemma}

\begin{proof}
BFK の対角 MF 公式
\cite[Theorem~5.39 and proof, pp.~82--84]{NFBFKKernels}の入力は、
$N$ 個の次数一変数、次数三の potential、および孤立臨界点という条件である。
本節の一般の $G$ はこれらを満たす。
$H^a\operatorname{RHom}(\mathbf1,\varrho^b)$ の恒等 sector は、
$a$ が偶数のとき $V_{b+3a/2}$、奇数のとき零である。
二つの非恒等 $\mu_3$ sector は原点だけを固定し、
それぞれが寄与する条件は
\[
                  a\equiv N\pmod2,\qquad 2b+3a=N
\]
である。特に境界の Calabi--Yau 次元は $k$ であり、
偶数 $k$ の parity を奇数 $k$ へ流用してはいない。
$b=2,a=-j<0$ とすると、恒等 sector は奇数 cohomological 次数では零、
偶数 cohomological 次数では負の Jacobian 次数となり、
非恒等 sector の必要条件は $4-3j=N\ge6$ となり不可能である。
これが最後の消滅を与える。

$U_r=\operatorname{RHom}(\mathbf1,J^r)$、
$W_r=\operatorname{RHom}(J,J^r)$、
$M_r=\operatorname{RHom}(K_\eta,J^r)$ と置く。
余単位三角から $U_r\to M_r\to W_r\to U_r[1]$ を得る。
$d=2-r-i$ に対する完全列の必要部分は
\[
 H^{d-1}W_r\longrightarrow H^dU_r\longrightarrow H^dM_r
       \longrightarrow H^dW_r
\]
である。その三つの既知項の恒等 sector は順に
\[
 \begin{array}{c|c|c}
 &\text{現れる parity}&\text{Jacobian 次数}\\\hline
 H^{d-1}W_r&i\text{ 偶数}&1-r-3i/2\\
 H^dU_r&i\text{ 偶数}&3-r-3i/2\\
 H^dW_r&i\text{ 奇数}&(5-2r-3i)/2
 \end{array}
\]
となる。両端は常に零であり、中項は $(r,i)=(2,0),(3,0)$ だけに残る。
非恒等 sector の必要条件はそれぞれ
\[
 2-2r-3i=N,\qquad 6-2r-3i=N,\qquad 5-2r-3i=N
\]
である。左辺は高々二なので全て不可能である。
従って $H^dU_r\to H^dM_r$ は同型となり、式\eqref{euk:eq:arity}を得る。

$U_r,W_r$ は固定境界の複体の底変更であり、cohomology は $R$ 上自由である。
同じ三角と両隣接群の消滅が全 $\eta$ について成立するため、
ここでは $M_r$ 自体を一定複体と同一視していない。
得た cohomology の同定は $\epsilon$ による前合成から来ており、底変更に自然である。
固定 counit の unit-insertion cube の total fiber は、
分裂を仮定せず $J^r$ となる。
補題\ref{cmo:lem:normalized-comonad-uniqueness}の opposite と Ind 完備化による適用確認を用い、
\cite[Theorem~3.3, Lemma~3.11]{CMOWillumsgaardAssociativity}の
正規化塔の非空 fiber の $i$ 次 homotopy 群を
$H^{2-r-i}M_r$ と同定すれば、残る主張が従う。
\end{proof}

\begin{lemma}[可換実族は必要な三段の差を実現する]
\label{euk:lem:commutative-reference}
一般の $G$ に対する可換三次族
\[
                    W=G+uQ+u^2L+u^3d
\]
の中心超平面制限を考える。固定した符号規約で、その twist 射は
$\eta^{\mathrm{cl}}=-[Q]\in V_2$ である。
従って二次式から $V_2$ への固定線型切断を選べば、任意の counit arrow を実現できる。
同じ $Q$ の下で $L\mapsto L+\delta L$ が与える二項の差は
$-\delta L\in V_1$ であり、同じ $Q,L$ の下で
$d\mapsto d+\delta d$ が与える三項の差は
$s_{\mathrm{cl}}\delta d$、$s_{\mathrm{cl}}\in\C^\times$ である。
これらは全係数上で成立する。
\end{lemma}

\begin{proof}
境界 MF の微分 $d_M$ は $d_M^2=G\operatorname{id}$ を満たす。
$W=0$ の環上では
\[
              d_M^2=-u(Q+uL+u^2d)\operatorname{id}
\]
なので、一次 Bockstein は $-Q$ の乗法である。
中心元の二項解消を有限 mutation へ通す
補題\ref{cmo:lem:spherical-bockstein}の比較により、これは残余の twist 射に等しい。
$\partial_jG$ の乗法は $\partial_jd_M$ によって自然に nullhomotopic であり、
BFK の恒等 sector 公式はその商を正確に $V_2$ と同定する。
したがって固定切断 $\sigma:V_2\to\C[\mathbf x]_2$ によって
$Q=-\sigma(\eta)$ を取れる。

$Q=0$ での二項の差は、$v=u^2$ に対する法線二列解消から計算できる。
補題\ref{cmo:lem:central-universal-parameters}の式
$D=\widetilde d+u(\partial_e-\widetilde\theta e)$、
$\xi=\partial_e+\widetilde\theta e$、$\xi^2=\widetilde\theta$ において、
可換側の $v$-Bockstein は $-\delta L$ の乗法である。
この解消は $G$ の特定の係数を使っておらず、
$J^2\simeq\varrho$ の周期を通して二項の差を与える。

$Q=L=0$ の三項の差は、中央一変数の rank 一 MF
\[
 d_0=\begin{pmatrix}0&u\\du^2&0\end{pmatrix},\qquad
 \theta=\begin{pmatrix}0&1\\-du&0\end{pmatrix},\qquad
 h=\begin{pmatrix}0&0\\-d&0\end{pmatrix}
\]
で計算する。
$\theta^2=d_0h+hd_0=-du\operatorname{id}$、
$\theta h+h\theta=-d\operatorname{id}$ なので、
正規化三項は $\pm d$ となる。
この法線因子を境界の非零単位へ tensor すれば、
補題\ref{cmo:lem:central-ternary-nonzero}と同じ実随伴の計算として
$V_0=\C$ の非零元を得る。

全係数への延長には重みを使う。
可換実随伴の構造は全多項式係数上にあり、前補題によって二項・三項の差は
$V_1$ と $\C$ に値を持つ大域多項式である。
重み $(Q,L,d)=(1,2,3)$ の下で、二項の差は重み二であり、
$\delta L=0$ で消えるので $\delta L$ の定数線型写像に限られる。
既に $Q=0$ でその写像は $-\operatorname{id}$ と分かった。
三項の差も $\delta d$ の非零定数倍に限られる。
三項を比較する前に二項を同定する path の空間は、
前補題の高次 homotopy 群の消滅により可縮である。
従って三項差の大域的定義に、局所的な原点選択や非多項式的な gauge は残らない。
\end{proof}

\subsubsection{全係数の三角変換と Morita 核}

\begin{proposition}[定数可逆行列による三角多項式変換]
\label{euk:prop:coefficients}
ある線型写像 $Q:\C^{3k}_{\mathbf t}\to\C[\mathbf x]_2$ と、
多項式
\begin{equation}\label{euk:eq:coefficients}
 \begin{split}
 L(\mathbf t,\alpha)&=\mathsf A\alpha+R_2(\mathbf t),\\
 C(\mathbf t,\alpha,c)&=\kappa c+R_{1,1}(\mathbf t,\alpha)+R_3(\mathbf t)
 \end{split}
\end{equation}
が存在する。ここで $\mathsf A:\C^{3k}_\alpha\to V_1$ は定数線型同型、
$\kappa\in\C^\times$、$R_2$ は $V_1$ に値を持つ斉次二次式、
$R_{1,1}$ は双線型式、$R_3$ は斉次三次式である。
非可換族と可換族 $G+uQ+u^2L+u^3C$ の実制限 comonad は、
同じ underlying counit arrow の上で同じ $A_3$ 成分に属する。
さらに
\[
 \Theta:(\mathbf t,\alpha,c)\longmapsto
     (\mathbf t,\ell=\operatorname{coeff}L,\gamma=C)
\]
は多項式同型である。
\end{proposition}

\begin{proof}
命題\ref{emc:prop:line-model}は全 $T$ 上の実随伴とその全構造を与え、
補題\ref{ebk:lem:boundary}は固定境界と rotation を可換側へ移す。
その twist 射の像を $\eta(\mathbf t)\in V_2\otimes T$ とする。
一次 Bockstein は $A/(u^2)$ のみから決まるため、$\alpha,c$ に依存しない。
法線の rescaling の重み一から、これは $\mathbf t$ に線型である。
ここでは二捩れ因子の中心平方の補正も $A/(u^2)$ に含めている。
前補題により $Q(\mathbf t)=-\sigma(\eta(\mathbf t))$ を取る。
二射の同じ cohomology 類を同定する homotopy は、
固定境界の mapping complex の $\C$-線型縮約を $T$ へ tensor して選べる。
法線の $\mathbf G_m$-作用による重み分解を使い、この縮約を重みを保つものとして取る。
負次数 Hom の消滅により、その path の選択空間は可縮である。
これで underlying counit arrow が全 $T$ 上で同定される。

まず実可換族 $G+uQ(\mathbf t)$ を二項構造の参照点に取る。
非可換側の実構造との差は、式\eqref{euk:eq:arity}の
自由加群 $V_1\otimes T$ の大域切断である。
可換側で $L$ を変える作用は全 $V_1$ を覆うため、
この差を消す $L(\mathbf t,\alpha)$ が多項式的に定まる。
重み二の単項式は $\alpha$ と $\mathbf t$ の二次積だけなので、
$L=\mathsf A\alpha+R_2(\mathbf t)$ という形になる。
$\mathbf t=0$ の二項は中心 $v=u^2$ 変形の Bockstein であり、
命題\ref{ebc:prop:constant-iso}により、その境界比較後の写像は同型である。
可換側の符号を合わせたものが $\mathsf A$ である。
共通の scalar であることを仮定せず、定数可逆行列として扱っている。

次に実可換族 $G+uQ(\mathbf t)+u^2L(\mathbf t,\alpha)$ を参照する。
二項構造は既に同定され、その同定の path の空間は可縮である。
従って三項構造の差は、自由な階数一加群 $T$ の大域切断として定まる。
可換側の $d$ による作用が非零定数倍なので、これを消す $C$ を多項式的に選べる。
重み三により $C$ は式\eqref{euk:eq:coefficients}の形になる。
$\kappa$ の非零性は $\mathbf t=\alpha=0$ で確認できる。
このとき source の曲率は $w_0+cu^3$ の分離和であり、
restriction--inflation の実解消の法線因子は前補題の rank 一 MF である。
従って source でも三項の $c$ 係数は非零の単位 class となる。
境界同値と指定周期の同定はこの一次元空間の同型なので、
可換側との比 $\kappa$ は非零である。
この検出は因子の同一性、PI 性、判別式の一致を用いない。

以上の構成では全係数上の二つの実構造の差を取っている。
各閉点で選んだ minimal model の貼合せや、三項 torsor の仮想的な原点は使っていない。
両隣接群の消滅による自然な同定が、各段の差と法線 rescaling の重みを保つ。
従って係数は全 $T$ 上の斉次多項式であり、局所化した有理関数ではない。

Morita 比較の存在に至る構成には、ここまで $\mathsf A$ の可逆性を必要としなかった。
可逆性と $\kappa\ne0$ を用いると、表示式から $\alpha$、続いて $c$ を
定数だけで割って多項式的に回復できる。これが $\Theta$ の同型性を示す。
$\mathbf t\mapsto\eta(\mathbf t)$ の単射性は、この結論にも不要である。
\end{proof}

\begin{theorem}[任意の固定楕円因子の全係数比較]
\label{euk:thm:universal-kernel}
式\eqref{euk:eq:coefficients}により
\[
 X_{\mathbf t,\alpha,c}
  =V\bigl(G+uQ(\mathbf t)+u^2L(\mathbf t,\alpha)
                    +u^3C(\mathbf t,\alpha,c)\bigr)\subset\mathbb P^N
\]
と定める。全係数環 $T$ 上で、非可換族とこの可換族の残余線束圏の間に
実 dg Morita 双加群と逆双加群が存在する。
それらは中心超平面制限を intertwine し、単位・余単位とともに任意の導来底変更に保たれる。
各複素係数で
\[
 \mathcal A_A\text{ が smooth}
       \quad\Longleftrightarrow\quad
 X_{\mathbf t,\alpha,c}\text{ が滑らか}
\]
である。可換族が滑らかな開集合上では比較核は相対積上 perfect となり、
各複素ファイバーの全 coherent 残余圏について
\[
          \Ku_{\mathrm{coh}}(A_{\mathbf t,\alpha,c})
             \simeq\Ku(X_{\mathbf t,\alpha,c})
\]
を得る。
\end{theorem}

\begin{proof}
前命題は、全 $T$ 上で counit arrow と実 $A_3$ 構造を同定する。
補題\ref{euk:lem:arity}により四項以降の非空 fiber は可縮である。
両側には最初から実随伴の全構造が存在するため、
比較 path は各段で持ち上がり、これらの可縮な選択の逆極限が
全 homotopy-coherent comonad の同値を与える。
二項・三項の $\pi_0$ の一致だけで止めているのではない。
また任意の仮想的な低項データに対する存在障害の消滅を仮定していない。

固定境界の生成対象を $G_{\mathcal B}$ とする。
命題\ref{emc:prop:line-model}の生成性により、$F^RG_{\mathcal B}$ は source の生成対象である。
実随伴の coKleisli 模型は
\[
 \operatorname{REnd}(F^RG_{\mathcal B})
   \simeq\operatorname{RHom}(K_\eta G_{\mathcal B},G_{\mathcal B}),
 \qquad g\star f=gK_\eta(f)\Delta_{G_{\mathcal B}}
\]
である。全 comonad 同値はこの積と高次構造を同定する。
その graph 双加群を両側の生成対象と tensor すれば、
全 $T$ 上の実 Morita 双加群と逆双加群を得る。
生成対象上の制限作用との互換性は全 comonad 比較から来ており、
有限 cone と直和因子へ延ばせば制限の intertwiner、単位・余単位を得る。
必要な核は固定 smooth proper 境界の相対 perfect 核であり、
mapping 複体、有限 mutation と選んだ homotopy は導来底変更に両立する。
従って構成全体が任意の導来係数変更に保たれる。

可換 smooth 開では、その残余線束圏は相対 smooth proper であり、
Morita 同値によって source の有限線束模型も同じ性質を持つ。
graph の underlying $T$-複体は perfect である。
両側の smooth な対角の有限解消を tensor すると、
graph は二側の相対積上も perfect となる。
周囲代数の noetherian Koszul AS-regular 性と正則中心三次は
命題\ref{emc:prop:line-model}の入力である。
有限例外列を戻し、主稿の命題\paperref{main}{crs:prop:residual-smooth}と
定理\paperref{main}{crs:thm:coherent}を適用すれば、全 coherent 圏への飽和を得る。

可換側の全 twist の線束圏は $\operatorname{Perf}(X_{\mathbf t,\alpha,c})$ である。
その残余が smooth なら有限例外列との貼合せで全線束圏も smooth となり、
同じ飽和判定により全 coherent 対象が perfect となる。
特異点があればその剰余体は hypersurface 局所環上で無限射影次元を持つので矛盾する。
逆向きは滑らかな超曲面の標準分解による。
全係数の Morita 同値が smoothness を保つことにより、表示した必要十分条件を得る。
特異係数の全 coherent 圏や、非正則な底上の全 coherent 圏の比較は結論していない。
\end{proof}

\subsubsection{判別行列と正規化の帰結}

\begin{corollary}[一般の楕円因子でも階数 $8^k$ の判別]
\label{euk:cor:discriminant}
$R=\C[\mathbf t,\alpha]$、$f=G+Q(\mathbf t)+L(\mathbf t,\alpha)$ と置く。
臨界代数と判別多項式を
\[
 \mathcal Q=R[\mathbf x]/(\partial f),\qquad
 D(\mathbf t,\alpha,c)
   =\det\bigl(C(\mathbf t,\alpha,c)\operatorname{id}_{\mathcal Q}+m_f\bigr)
\]
と定める。$\mathcal Q$ は階数 $\mu=2^N=8^k$ の有限自由 $R$-加群であり、
\[
                    [f]=\left[\frac{Q+2L}{3}\right]\in\mathcal Q
\]
である。$D$ の $c$ に関する次数は $\mu$、最高係数は $\kappa^\mu$ であり、
\[
 X_{\mathbf t,\alpha,c}\text{ が滑らか}
    \quad\Longleftrightarrow\quad D(\mathbf t,\alpha,c)\ne0
    \quad\Longleftrightarrow\quad\mathcal A_A\text{ が smooth}
\]
となる。任意の固定した $\mathbf t$ について、一般の $\alpha$ では
正確に $8^k$ 個の相異なる悪い中心係数があり、各可換標的は唯一の通常二重点を持つ。
\end{corollary}

\begin{proof}
$\partial G$ は原点だけを共通零点とする二次式の正則列である。
active 濾過の Koszul 複体の随伴はこの正則列の複体なので、
$\operatorname{gr}\mathcal Q=\operatorname{Jac}(G)\otimes R$ となる。
固定した Jacobian 基底の持上げは $R$-基底を与える。
Hilbert 級数 $(1+t)^N$ より階数は $2^N$ である。
Euler の式 $3G+2Q+L=0$ を臨界代数で使えば表示した $[f]$ を得る。
$u=0$ には特異点がないので、特異点の存在は $u=1$ で
$\partial f=0$、$C+f=0$ を満たす点の存在と同値である。
有限代数の元が可逆であることと乗法行列式の非零性は同値だから、判別条件が従う。
特殊化した臨界代数を $\mathcal Q_{\mathbf t,\alpha}=\mathcal Q\otimes_R\C$ とすると、
より正確に
\[
 D=\prod_{p\in(\operatorname{Spec}\mathcal Q_{\mathbf t,\alpha})(\C)}
                  (C+f(p))^{\operatorname{length}(\mathcal Q_{\mathbf t,\alpha})_p}
\]
と読む。ここで多重臨界値と非被約な臨界点を除いていない。

固定した $\mathbf t$ では $Q$ は固定され、$\alpha\mapsto L$ は
命題\ref{euk:prop:coefficients}により affine 線型同型である。
補題\ref{ldg:lem:fixed-quadratic-morsification}の証明は、
最高次の $\partial G$ が正則列であることだけを用いるので、この $G$ に適用できる。
具体的には $-\nabla(G+Q)$ は次数 $\mu$ の有限平坦射であり、
標数零により一般のファイバーは \`etale である。
相異なる臨界点対での臨界値差 $h$ は
$dh=\sum_j(x_j-y_j)d\ell_j$ を満たすので、どの成分でも恒等零にはならない。
その零集合の有限像を除けば全臨界値が相異なる。
$c\mapsto C$ の係数は非零の $\kappa$ なので、各値は一つの悪い $c$ に対応する。
非退化 Hessian と形式的 Morse 補題が、そこでの唯一の通常二重点を与える。
\end{proof}

\begin{corollary}[判別超曲面の既約性と affine 正規化]
\label{euk:cor:normalization}
$D$ は全係数環 $\C[\mathbf t,\alpha,c]$ 上で既約であり、
任意の固定した $\mathbf t$ について $\C[\alpha,c]$ 上でも既約である。
全係数の判別超曲面の正規化空間は
$\mathbb A^{3k}_{\mathbf t}\times\mathbb A^N_{\mathbf x}$ であり、
固定した $\mathbf t$ では $\mathbb A^N_{\mathbf x}$ である。

$\Theta$ の座標 $(\mathbf t,\ell,\gamma)$ を使い、
$H_{\mathbf t}=G+Q(\mathbf t)$ と書けば、正規化射は
\[
 (\mathbf t,\mathbf x)\longmapsto
 \bigl(\mathbf t,-\nabla H_{\mathbf t}(\mathbf x),
       \mathbf x\cdot\nabla H_{\mathbf t}(\mathbf x)-H_{\mathbf t}(\mathbf x)\bigr)
\]
である。複素点として、判別超曲面の特異点は、
値 $-C$ を取る異なる臨界点がある係数、または値 $-C$ の臨界点で Hessian が退化する係数に
正確に対応する。これらは固定 $\mathbf t$ の判別超曲面にも成立する。
\end{corollary}

\begin{proof}
三角多項式同型 $\Theta$ へ移すと、臨界スキームは
$\ell=-\nabla H_{\mathbf t}$ によって affine 空間
$\mathbb A^{3k}_{\mathbf t}\times\mathbb A^N_{\mathbf x}$ に等しい。
これは整かつ正規であり、$(\mathbf t,\ell)$ 上有限平坦である。
一般係数体での norm 多項式は $-f$ の最小多項式の冪であるが、
前系の臨界値の相異性がその指数を一にする。
monic 性と Gauss の補題から既約性を得て、
有限・双有理な表示射が正規化となる。
この議論は各固定 $\mathbf t$ にも成立する。
微分の核は $\{(0,v):\operatorname{Hess}H_{\mathbf t}\,v=0\}$ である。
従って定理\ref{ldg:thm:normalization}の局所環の証明、すなわち
一点ファイバーかつ余接写像が全射であるときの二段の Nakayama の議論を適用でき、
特異点の必要十分条件を得る。
\end{proof}

\begin{remark}[固定因子の比較の範囲]
\label{euk:rem:scope}
本節の $\kappa$ と係数行列は、固定した因子、方程式と境界比較の規約に依存する。
全因子を $S(1,-1,-1)$ とし、前節の方程式と二項の正規化を用いる場合には、
系\ref{mcs:cor:universal-scale}の $\kappa_k^2=27/28$ を回復する。
一般の異なる楕円因子について、この数値を流用してはいない。
$Q,R_2,R_{1,1},R_3$ の全数値係数も本節では求めていない。

全係数と各固定 $\mathbf t$ での判別超曲面の被約性は、
固定した $(\mathbf t,\alpha)$ での一変数多項式の被約性を意味しない。
特殊化後には臨界値の衝突と臨界代数の非被約性を許す。
また factor の楕円点自体を動かす相対族、全 intrinsic 対象の表示被覆、
核による全中心作用・rotation・周期標識の同時保存は結論していない。
本節が与えるのは、指定した全中心持上げ係数に対する残余線束圏の実比較と、
その可換 smooth 開における相対 perfect 核および全 coherent Ku 比較である。
\end{remark}

